\clearpage
\section{Semantic soundness of structural derivations and resource renaming}
\label{sec:copying-semantic-renaming}
The full BPerm robustness theorem now gives an operational meaning to every
structural BPerm derivation between admissible open programs. Its conclusion
quantifies over all the original binding-free closing graph contexts, including
contexts whose selected occurrence lies beneath arbitrarily many boxes.
Successful runs correspond with exactly the same number of machine transitions.
Applying this result to the existing full-syntax renaming derivation proves
semantic invariance under bijective renaming of both variables and atoms, with
the external resource permutation explicitly compensated.

This completes the semantic interpretation of the structural bijective-renaming
derivations in Appendix~\ref{app:remaining-work}. The terms, distributors, boxes,
padding and Figure~16 translation retain their existing definitions. In
particular, a permutation of free-resource positions is an explicit wiring
operation outside the translated graph. No claim identifies a nontrivial box
permutation with an ordered graph isomorphism. Non-bijective renaming, semantic
substitution and weakening, the remaining template proofs and the equivalence
between the two different terms of Proposition~12.2 are still separate tasks.

The proof has three stages. First, interpret a structural derivation by a finite
chain of original recursive filling certificates. Second, lift each certificate
through the chosen addressed graph context. Third, apply full BPerm robustness
to each closed link and compose the resulting exact machine equivalences.
Only the machine equivalences are composed transitively; no transitivity law
is added to the original filling relation.

\subsection{Primitive-label properties along recursive fillings}
Let $P$ be any predicate on primitive labels. Assume that each pair in a
pre-template preserves the assertion that every primitive label satisfies $P$.
Then every original recursive paired filling preserves that assertion too.
The surrounding context itself need not satisfy $P$: for example, its hole
labels may fail $P$ even though filling those holes produces hole-free graphs.

The proof is induction on the box allowance. At allowance zero, each component
is either a retained corolla, whose two labels are equal, or a genuine filled
hole, where the template hypothesis applies. At a successor allowance there
is also a boxed component. Read the label property inside the source box,
apply the induction hypothesis to its body filling, and rebuild the target
box. Every expanded edge is an original component edge, so these component
equivalences assemble into the result for the expanded graph.

\lean{Copying/DerivationFillings}{14}{57}{DeepFilling.allLabels_iff}

A filling witness also contains recursive isomorphisms from the expansions
to the actual endpoints. Primitive-label invariance under those isomorphisms
transports the preceding equivalence to the endpoints. BPerm preserves every
primitive-label predicate, since it changes box interfaces and the associated
wiring without changing the primitive labels in the bodies. Its specialization
therefore needs no additional predicate-preservation premise.

\lean{Copying/DerivationFillings}{59}{65}{FillingWitness.allLabels_iff}
\lean{Copying/DerivationFillings}{69}{72}{filling_allLabels}

\Needspace{11\baselineskip}
\subsection{Finite chains of original certificates}
\termdef{finite-bperm-filling-chain}
A \emph{finite BPerm filling chain} from $G$ to $H$, at a common interface
and allowance, consists of a natural number $N$, graphs $G_i$ for all
$i\in\mathbb N$, and original filling witnesses between every $G_i$ and
$G_{i+1}$, with
\[
 G_0=G,\qquad G_i=H\quad\text{whenever }N\leq i.
\]
The constant tail lets a completed component remain unchanged while another
component continues its derivation. The number $N$ bounds the structural
certificate chain; it is not a count of machine transitions.

\lean{Copying/DerivationFillings}{76}{81}{FillingChain}

Every $G_i$ is focus-free because it is the left endpoint of an original
filling witness. Induction on $i$ also proves that $G$ and $G_i$ satisfy
exactly the same universal primitive-label properties: the base case uses
$G_0=G$, and the successor case composes the equivalence for the next witness.
Taking $P$ to exclude holes will later make every intermediate closed graph
an admissible machine input.

\lean{Copying/DerivationFillings}{85}{86}{FillingChain.focusFree}
\lean{Copying/DerivationFillings}{88}{92}{FillingChain.allLabels}

The reflexive chain has bound zero and constantly repeats its focus-free
endpoint, using the original unchanged-edge corollas for every certificate.
A single original witness gives a chain of bound one: use that witness at
index zero and repeat its right endpoint thereafter. Its focus freedom is a
consequence of the supplied witness, so the tail needs no new assumption.

\lean{Copying/DerivationFillings}{95}{100}{FillingChain.refl}
\lean{Copying/DerivationFillings}{103}{112}{FillingChain.single}

Concatenating chains of bounds $N$ and $M$ gives bound $N+M$. Before $N$ use
the first chain, and from $N$ onward use index $i-N$ of the second. At the
junction, the first chain's endpoint equation and the second chain's initial
equation identify the common graph. The proof also handles $N=0$ explicitly.
Thus each link remains one of the supplied original certificates.

\lean{Copying/DerivationFillings}{116}{135}{FillingChain.trans}

Reversal uses graph $G_{N-i}$, where subtraction is truncated at zero.
Before the bound, reverse the corresponding original filling by the previously
proved \texttt{reverseFilling} construction. At and after the bound the graph
is constantly $G_0$, so use its reflexive certificate. This construction uses
proved reversal of whole filling witnesses and does not assume symmetry of
the directed raw BPerm pair.

\lean{Copying/DerivationFillings}{138}{150}{FillingChain.symm}

\termdef{bperm-chain-synchronisation}
\emph{Chain synchronization} applies a graph constructor pointwise to chains
whose completed tails repeat their endpoints. For serial composition and
binary tensor, use the larger of the two bounds. By that point both component
chains are constant, and the existing serial or tensor operation on original
filling witnesses supplies every link. Their start and finish equations follow
by substituting the component endpoint equations.

\lean{Copying/DerivationFillings}{153}{160}{FillingChain.serial}
\lean{Copying/DerivationFillings}{162}{169}{FillingChain.tensor}

For a finite parallel family, sum the component bounds over the supplied
enumeration. Each component bound is at most this sum. The elementary bound
lemma is induction on the enumeration list: an occurrence is either its head,
where nonnegativity suffices, or in its tail, where the induction hypothesis
applies. Repetitions in the enumeration are harmless. The parallel chain uses
the existing common boundary assembly, and advances all components at the
same index. Completeness of the enumeration proves that every component has
reached its endpoint by the chosen sum; the parallel filling construction
supplies each link.

\lean{Copying/DerivationFillings}{173}{182}{FillingChain.length_le_sum}
\lean{Copying/DerivationFillings}{186}{199}{FillingChain.parallel}

Boxing a chain preserves its bound and boxes each graph with the same binding
metadata; the original boxed-witness construction supplies its links.
Transporting equal input and output interface lists similarly transports
every graph and certificate without changing the bound. Raising the
representation allowance applies the existing raising operation to every
graph, and uses BPerm's proved padding closure to raise each original witness.
These three constructions inherit their endpoint equations pointwise.

\lean{Copying/DerivationFillings}{201}{209}{FillingChain.boxed}
\lean{Copying/DerivationFillings}{211}{218}{FillingChain.interfaces}
\lean{Copying/DerivationFillings}{220}{226}{FillingChain.raise}

Transport along an equality of allowance indices is equality elimination:
after identifying the indices it is the original chain. Padding by an
arbitrary number of levels is recursion on that number, with the original
chain at zero and one more raising operation at the successor. These operations
change representation allowance without adding structural chain links.

\lean{Copying/DerivationFillings}{228}{230}{FillingChain.depthCast}
\lean{Copying/DerivationFillings}{232}{235}{FillingChain.pad}

Every structural BPerm derivation whose source is focus-free now has such a
finite chain. Induct on the derivation. An isomorphism changes one endpoint
of a reflexive original witness; a BPerm rule supplies a genuine filled-hole
witness. Symmetry and transitivity use the chain operations above. Serial,
tensor and finite parallel derivations use their synchronized chains; boxing,
interface transport and raising use their pointwise constructions. Where the
induction moves to components, extract their focus freedom from that of the
assembled source. For symmetry and the second part of transitivity, use the
already proved focus-freedom invariance of structural derivations. The final
definition selects one of the chains whose existence has just been proved.
Classical choice selects certificate data here; it supplies no semantic premise.

\lean{Copying/DerivationFillings}{241}{274}{Derivation.fillingChain_exists}
\lean{Copying/DerivationFillings}{276}{278}{Derivation.fillingChain}

\subsection{Lifting certificates through addressed graph contexts}
\termdef{addressed-filling-transport}
\emph{Addressed filling transport} places both endpoints of an original
recursive filling witness into the same selected occurrence of a focus-free
host. It works for any pre-template and any well-typed addressed occurrence;
the occurrence may lie inside nested original boxes.

For a shallow edge $e$ of a host $C$, first form the simultaneous expansion
that inserts the left or right endpoint at $e$ and the original corolla at
every other edge. Use the supplied witness at $e$. For every other edge use
the retained-corolla certificates; focus freedom follows by reading that
edge's original label in $C$. The expansion operation on filling witnesses
assembles these components. The proved single-expansion isomorphisms then
identify the two endpoints with the actual single-edge plugs.

\lean{Copying/AddressedWitness}{11}{27}{FillingWitness.singleExpansion}
\lean{Copying/AddressedWitness}{30}{34}{FillingWitness.plug}

For an enclosing box, use the body witness to obtain witnesses between the
two rebuilt box corollas. Keep its bound-variable count, free-variable count,
free-atom count and outer incidences fixed. Insert that component into the
host expansion, then use the box-expansion isomorphisms to identify the
endpoints with the actual box replacements.

\lean{Copying/AddressedWitness}{37}{44}{FillingWitness.replaceBox}

Induction on the address now proves the general result. A shallow address
uses single-edge plugging. For an address inside a box, extract focus freedom
of the original body from the host, apply the induction hypothesis there,
and rebuild the enclosing box by the preceding construction. No new filling
constructor is required. Specializing to a closing observation context uses
the focus-freedom component of its original binding-free hypothesis.

\lean{Copying/AddressedWitness}{48}{58}{FillingWitness.plugAt}
\lean{Copying/AddressedWitness}{61}{63}{FillingWitness.closing}

Finally apply this operation to every graph and every link of a filling chain,
using the same host and address throughout. Its bound is unchanged, and
its endpoint equations follow by applying the plug operation to the old
endpoint equations. This yields a chain between the actual addressed fillings
of the original two endpoints.

\lean{Copying/ContextualDerivation}{11}{18}{FillingChain.plugAt}

\subsection{Exact execution and the full observation class}
\termdef{structural-bperm-semantic-soundness}
The \emph{semantic soundness of a structural BPerm derivation} states that
admissible open endpoints are exactly equivalent in the original copying
machine under every allowed binding-free closing graph context. Here
``exactly'' refers to successful machine-transition counts, not to equality
of the endpoint graphs or their final returned graph representations.

First consider a chain at closed interface $[\mathsf{term}]\to[]$, whose
source contains no holes. Label preservation proves that all intermediate
graphs contain no holes. Each adjacent original filling witness therefore
satisfies the premises of the full BPerm exact-execution theorem from
Section~\ref{sec:copying-full-bperm}. Induction on the chain index composes the
two directions of these machine refinements. At index zero use reflexivity;
at a successor compose the previous comparison with the next link. At the
chain bound, the finish equation identifies the graph with the target.

\lean{Copying/DerivationExecution}{13}{29}{FillingChain.exact_equiv}

Unfolding exact refinement gives a useful count-by-count statement. For each
$n$, a successful run of exactly $n$ transitions exists from the source if
and only if such a run exists from the target. In either direction the
refinement provides a target count $m$ and its exact-cost equality with $n$;
substitute that equality to obtain the stated count. Applying these two results
to the chain extracted from a focus-free structural derivation gives their
closed-derivation forms. Hole freedom is required only at the source; it is
proved for the rest of the chain.

\lean{Copying/DerivationExecution}{33}{46}{FillingChain.termination_count}
\lean{Copying/DerivationExecution}{50}{54}{Derivation.exact_equiv}
\lean{Copying/DerivationExecution}{56}{60}{Derivation.termination_count}

Now let $G,H$ be open programs and $s$ a structural derivation between their
graphs. Given any original closing context $C$, extract a chain for $s$ and
lift it at $C$'s designated occurrence. Binding freedom provides the host's
focus freedom. The context's unique-hole inventory and the source program's
hole freedom prove that $C[G]$ has no holes. Apply the closed-chain theorem
to conclude exact equivalence of the initial machine states of $C[G]$ and
$C[H]$. Quantification over $C$ is precisely the existing definition of
contextual equivalence; there is no restriction to translated term contexts
and no bound on the surrounding context's nesting beyond its own finite
representation.

\lean{Copying/ContextualDerivation}{22}{27}{Derivation.contextual_exact}

The contextual count theorem applies the same lifted-chain argument to the
explicit successful-run formula. Forgetting costs then proves termination-based
contextual equivalence: from a successful observation, take its finite count,
transfer that count, and retain the resulting successful observation. The
reverse implication is identical with the endpoints exchanged. These results
assume neither endpoint nor callback terminates. They characterize successful
termination, and do not separately identify divergence and stuckness.

\lean{Copying/ContextualDerivation}{31}{36}{Derivation.contextual_count}
\lean{Copying/ContextualDerivation}{39}{49}{Derivation.contextual_termination}

\subsection{Semantic renaming with explicit boundary compensation}
\termdef{compensated-semantic-renaming}
Let $t$ have $k$ free-variable positions and $h$ atom positions. Let $\rho$
and $\alpha$ be bijections of these two finite sets, and write $p$ for the
typed output permutation combining them. Write $\lbrack\!\lbrack t\rbrack\!\rbrack_d$
for the original translation padded to an allowance $d$ at least the syntactic
depth of $t$. The \emph{compensated renamed program} is
\[
 \mathsf{perm}_{p^{-1}}
   \bigl(\lbrack\!\lbrack \mathsf{rename}_{\rho,\alpha}(t)\rbrack\!\rbrack_d\bigr).
\]
The inverse outer permutation reconnects each renamed resource to its original
context position. The two sorts are kept separate, and the variable renaming
under a binder fixes the newly bound variable. Compensation is an actual
external output-wiring operation; it does not alter the definition of term
translation or license permutation of the bound prefix.

The existing structural renaming certificate has a derivation followed by a
boundary action, potentially presented at heterogeneous allowance indices.
Pad that certificate first to one common sufficient allowance. Its existing
derivation theorem then yields the desired compensated endpoints directly;
the now-equal allowance cast simplifies away. Applying this construction to
the full-syntax translation renaming certificate gives a structural derivation
between $\lbrack\!\lbrack t\rbrack\!\rbrack_d$ and the displayed compensated graph.
The previously proved preservation of syntactic depth supplies the same
allowance bound for the renamed term.

\lean{Copying/SemanticRenamingPadding}{11}{20}{BoxPermutation.Renaming.paddedDerivation}
\lean{Copying/SemanticRenamingPadding}{24}{31}{Translation.padded_term_rename_derivation}

Permuting an admissible open program changes only its external output
positions. Its primitive labels, and hence its proofs of focus freedom and
hole freedom, are unchanged. The definition of \texttt{renamedProgram}
applies this operation with $p^{-1}$ to the unchanged renamed translation
at the selected sufficient allowance.

\lean{Copying/SemanticRenaming}{12}{16}{OpenProgram.permute}
\lean{Copying/SemanticRenaming}{20}{25}{renamedProgram}

Apply contextual soundness of structural derivations to the padded renaming
derivation. It proves, for every original term and every sufficient $d$,
\[
 \lbrack\!\lbrack t\rbrack\!\rbrack_d
 \simeq_{\mathrm{exact}}
 \mathsf{perm}_{p^{-1}}
   \bigl(\lbrack\!\lbrack \mathsf{rename}_{\rho,\alpha}(t)\rbrack\!\rbrack_d\bigr).
\]
The proof uses all constructors of the existing term syntax through its
already established structural renaming theorem, including nested lambdas,
application, every unary operation and every binary operation. All original
closing graph contexts are quantified, so the statement is not limited to
closed terms or to contexts with terminating callbacks.

\lean{Copying/SemanticRenaming}{29}{34}{Translation.term_rename_contextual_exact}

The count-by-count form and the cost-forgetting form are direct applications
of the corresponding contextual derivation theorems to the same derivation.
In the first, choose an arbitrary closing context and count $n$; the successful
run formula holds in both directions at that exact count. In the second,
forget the count while retaining successful termination. The named-environment
form instantiates $k,h$ with the lengths of the given distinct variable and
atom lists and applies the same exact theorem. Neither list must be empty,
and no extra hypothesis is placed on the translated term.

\lean{Copying/SemanticRenaming}{38}{47}{Translation.term_rename_contextual_count}
\lean{Copying/SemanticRenaming}{51}{56}{Translation.term_rename_contextual_termination}
\lean{Copying/SemanticRenaming}{60}{66}{Translation.named_rename_contextual_exact}

\subsection{Regression instances and the Proposition 12.2 boundary}
The first regression uses the existing term with two nested lambdas around
an application of a free variable to an atom. At allowance two, simultaneously
swap the two ambient variable positions and the two atom positions, and apply
the exact semantic renaming theorem. The second places that term beside the
unit constant in an application, and compares at allowance three. Its branches
have original allowances zero and two, so it exercises the original maximum
padding inside application and an additional outer padding level. Both
proofs are instances of the general theorem, with the concrete depth bounds
checked by computation.

\lean{Copying/SemanticRenamingExamples}{12}{17}{nestedRenaming_contextual_exact}
\lean{Copying/SemanticRenamingExamples}{21}{26}{unequalDepthRenaming_contextual_exact}

The count regression chooses any closing graph context for that latter
program and any natural number $n$. Its conclusion states the equivalence
of the two explicit final-run predicates at count $n$, by the general
contextual count theorem. The context's own structure, effects and nesting
remain universally quantified subject to its original closing-context laws.

\lean{Copying/SemanticRenamingExamples}{30}{40}{unequalDepthRenaming_contextual_count}

Finally, apply exact semantic renaming separately to the unchanged two terms
of Proposition~12.2. Their syntactic definitions use no ambient variable or
atom, so renaming those ambient positions reduces to the same original term.
The resulting statements compare each translated term with its own externally
permuted graph, for arbitrary $k,h$, bijections and sufficient allowance.
They retain the original let/ref/lambda/sequence/dereference expression on
the left and the original lambda/sequence/constant expression on the right.
The named private-cell instance specializes its arbitrary scope sizes to the
lengths of distinct-list environments.

\lean{Copying/SemanticRenamingExamples}{44}{51}{privateCell_contextualRenaming}
\lean{Copying/SemanticRenamingExamples}{55}{62}{constantResult_contextualRenaming}
\lean{Copying/SemanticRenamingExamples}{66}{74}{privateCell_named_contextualRenaming}

These last results concern each term's resource renaming independently.
They do not identify the two different terms with each other. Establishing
that equivalence still requires the remaining template obligations and the
Figure~22 chain, followed by the semantic weakening argument for the final
general-environment proposition.
